\section{The macroscopic limit for positive Rayleigh--Jeans data on a smooth \texorpdfstring{\mbox{Euler}}{Euler} interval}
\label{hyd:section}

We consider the four-wave collision operator $\mathcal C$ defined in the
preceding section, with spatial domain $\T^3=(\R/2\pi\mathbb Z)^3$ and
velocity $v(k)=2k$. Writing $\varepsilon=4\pi/c$ for the equivalent small
parameter, set
\[
 \Psi(k)=(1,k^{(1)},k^{(2)},k^{(3)},|k|^2)^{\mathsf T},\qquad
 \mathcal T=\partial_t+v(k)\cdot\nabla_X,
 \qquad c=\frac{4\pi}{\varepsilon}.
\]
All momentum integrals are over $D=[-R,R]^3$. We prove that, for positive
local Rayleigh--Jeans initial data, the strong-collision equation has a
solution throughout any prescribed interval of smooth positive Euler
existence, and we obtain convergence rates for both the macroscopic
parameters and the distribution. The collision frequency vanishes at the
corners, so the proof treats both the usual dissipative region and the
corner layers that retain a free-transport trace. Here $v_j(k)=2k^{(j)}$;
the subscript in $k_i$ is reserved for the $i$th collision leg.
The proof is organized around the five moments of the same evolving
distribution. The exact entropy decomposition first identifies the
moment-matched Rayleigh--Jeans equilibrium and the microscopic deviation;
a third-order energy with a nonlinear corrector and a corner comparison
then control the original initial-value problem. The main result is
Theorem~\ref{hyd:main}: the distribution error is $O(c^{-3/4})$, whereas
the second-order Sobolev error of the macroscopic parameters is $O(c^{-1})$.
We first construct a local solution and the nonlinear corrector, establish
the joint energy and corner positivity estimates, and then use a
first-exit argument to reach the prescribed Euler interval.

\subsection{Five-moment parameters and the symmetric hyperbolic system}
\label{hyd:euler-section}

Define the positive parameter domain and its associated distributions by
\begin{equation}\label{hyd:thermodomain}
 \mathcal O=\{\theta\in\R^5:\min_{k\in D}\theta\cdot\Psi(k)>0\},
 \qquad N_\theta(k)=\frac1{\theta\cdot\Psi(k)}.
\end{equation}
The set $\mathcal O$ is open and convex. Set
\begin{align}
 U(\theta)&=\int_D\Psi N_\theta\,\dd k,
 &\mathcal F_j(\theta)&=\int_Dv_j\Psi N_\theta\,\dd k,
 \label{hyd:moment-map}\\
 M(\theta)&=\int_DN_\theta^2\Psi\Psi^{\mathsf T}\,\dd k,
 &J_j(\theta)&=\int_Dv_jN_\theta^2\Psi\Psi^{\mathsf T}\,\dd k.
 \label{hyd:matrices}
\end{align}
Direct differentiation gives $DU=-M$ and $D\mathcal F_j=-J_j$.
Every nonzero coefficient vector defines a nonzero polynomial
$b\cdot\Psi$, so $M$ is positive definite; each $J_j$ is real symmetric.
For $\theta,\vartheta\in\mathcal O$,
\begin{equation}\label{hyd:moment-injective}
 (\theta-\vartheta)\cdot[U(\theta)-U(\vartheta)]
 =-\int_D\frac{[(\theta-\vartheta)\cdot\Psi]^2}
 {(\theta\cdot\Psi)(\vartheta\cdot\Psi)}\,\dd k.
\end{equation}
Thus $U$ is injective, and the inverse function theorem makes it a smooth
diffeomorphism from $\mathcal O$ onto its open image $U(\mathcal O)$.
Parameters will be recovered from moments only within this image.

\begin{proposition}[Exact entropy decomposition under five-moment matching]\label{hyd:entropy-geometry}
 Fix a spatial point. Let $f>0$ almost everywhere, $f\in L^1(D)$,
 $N=N_\theta$, and $\bar N=N_{\bar\theta}$, where
 $\theta,\bar\theta\in\mathcal O$.
 Write $U_f=\int_D\Psi f\dd k$ and define the relative entropy, allowing
 the value $+\infty$, by
 \[
 \mathscr H(f\mid g)=\int_D\left(\frac f g-1-\log\frac f g\right)\dd k.
 \]
 Then the three-point identity is
 \begin{equation}\label{hyd:entropy-threepoint}
 \mathscr H(f\mid\bar N)
 =\mathscr H(f\mid N)+\mathscr H(N\mid\bar N)
       +(\bar\theta-\theta)\cdot[U_f-U(\theta)].
 \end{equation}
 In particular, if $U_f=U(\theta)$, then
 \begin{equation}\label{hyd:entropy-pythagoras}
 \mathscr H(f\mid\bar N)
       =\mathscr H(f\mid N)+\mathscr H(N\mid\bar N).
 \end{equation}
 Among positive distributions with these five moments, $N$ uniquely
 minimizes both $\mathscr H(f\mid\bar N)$ and
 $-\int_D\log f\dd k$, up to almost-everywhere equality.

 The parameter potential
 $\mathcal P(\theta)=-\int_D\log(\theta\cdot\Psi)\dd k$ and the entropy
 in moment coordinates
 $\mathcal E(U(\theta))=-\int_D\log N_\theta\dd k$ satisfy
 \begin{equation}\label{hyd:entropy-hessians}
 D_\theta\mathcal P=-U,\qquad D_\theta^2\mathcal P=M,
 \qquad \nabla_U\mathcal E=-\theta,\qquad D_U^2\mathcal E=M^{-1}.
 \end{equation}
 On every fixed compact convex set $\Theta\Subset\mathcal O$, all
 $\theta,\bar\theta\in\Theta$ satisfy
 \begin{equation}\label{hyd:entropy-comparison}
 \mathscr H(N_\theta\mid N_{\bar\theta})
 \asymp_\Theta|\theta-\bar\theta|^2
 \asymp_\Theta|U(\theta)-U(\bar\theta)|^2.
 \end{equation}
\end{proposition}

\begin{proof}
 Expanding the three entropy densities pointwise leaves
 $(f-N)(\bar N^{-1}-N^{-1})$. Both reciprocals are linear combinations
 of $\Psi$, so integration gives \eqref{hyd:entropy-threepoint}.
 The cross term is absolutely integrable. The two relative entropies
 involving $f$ are finite if and only if $\log^-f\in L^1(D)$: the positive
 logarithmic part is controlled by $\log^+f\le f$, and both RJ states have
 positive upper and lower bounds. If the negative logarithmic part is
 not integrable, both sides of \eqref{hyd:entropy-threepoint} equal
 $+\infty$; the additive identity involves no subtraction
 $\infty-\infty$. The solutions with positive two-sided bounds considered
 here satisfy the finiteness condition. Moment matching makes the cross
 term vanish. The inequality $x-1-\log x\ge0$, with equality only at $x=1$,
 proves the first unique-minimizer assertion. Moreover,
 $\int_D(f/N-1)\dd k=\theta\cdot[U_f-U(\theta)]=0$, so
 $-\int\log f+\int\log N=\mathscr H(f\mid N)$, proving the second.

 Direct differentiation of $\mathcal P$ gives the first two derivative
 formulas; $DU=-M$ and $M\theta=U$ give the other two. With
 $\delta\theta=\bar\theta-\theta$,
 \[
 \begin{aligned}
 \mathscr H(N_\theta\mid N_{\bar\theta})
 &=\mathcal P(\bar\theta)-\mathcal P(\theta)
                   -D\mathcal P(\theta)\cdot\delta\theta\\
 &=\int_0^1(1-s)\,\delta\theta^{\mathsf T}
         M(\theta+s\delta\theta)\delta\theta\,\dd s.
 \end{aligned}
 \]
 Uniform positive definiteness of $M$ on $\Theta$ gives the quadratic
 entropy comparison. Integrating $DU=-M$ along the same segment and
 pairing with $\delta\theta$ gives both bounds for the moment difference.
\end{proof}

A simple matching example is $N=1$. Choose a real $\zeta$ with
$|\zeta|R^2<1$ and set
\[
 f_\zeta(k)=1+\zeta k^{(1)}k^{(2)}.
\]
This distribution is strictly positive on $D$, and
\[
 \int_D\Psi(f_\zeta-1)\,\dd k
       =\zeta\int_D\Psi k^{(1)}k^{(2)}\,\dd k=0.
\]
Each component of the last integrand is odd in at least one coordinate.
Its moment-matched equilibrium is therefore exactly $N=1$, whereas
$\mathscr H(f_\zeta\mid1)>0$ when $\zeta\ne0$: equality of all five
moments still permits a nonzero microscopic deviation.

The preceding identity is the concrete three-point formula for the convex
entropy under five-moment constraints. It holds at each $(t,X)$ for the
same $f,N,\bar N$, hence also after spatial integration; the collision
term and its four-leg dissipation continue to be computed from the
original distribution $f$. Its infinitesimal form explains the
orthogonal energy used below. At $f=N$, suppose that bounded variations
$\delta f$ and $\delta N=-N^2\Psi\cdot\delta\theta$ have matching five
moments. Set $h=(\delta f-\delta N)/N$ and $\kappa_N=N\Psi$; then
$\int_D\kappa_Nh\dd k=0$. Since
$(-\log)^{\prime\prime}(N)=N^{-2}$, the second variation of the entropy is
\[
 \left\|\frac{\delta f}{N}\right\|_2^2
 =\|h\|_2^2+\delta\theta^{\mathsf T}M\delta\theta.
\]
Thus the zeroth-order orthogonality is compatible with the entropy
Hessian. The third-order differentiated energy with a corrector used
below is not a spatial derivative of the total entropy: its control
requires the full collision and corner-layer arguments. In particular,
zero initial microscopic entropy need not imply zero initial error
relative to the corrected approximation.

The five-moment conservation laws are equivalent to
\begin{equation}\label{hyd:euler}
 M(\theta)\theta_t+\sum_{j=1}^3J_j(\theta)\partial_j\theta=0,
 \qquad
 \theta_t+\sum_j A_j(\theta)\partial_j\theta=0,
 \qquad A_j=M^{-1}J_j.
\end{equation}
The positive matrix $M$ is a symmetrizer. Write
$J(\xi)=\sum_j\xi_jJ_j$. For every $\xi\in\R^3$,
$M^{-1/2}J(\xi)M^{-1/2}$ is real symmetric; repeated eigenvalues cause
no difficulty in the following construction. This local construction
belongs to the theory of quasilinear symmetric hyperbolic
systems~\cite{Kato1975}; we give the energy and approximation arguments
needed here.

\begin{proposition}[Local positive Euler solutions]\label{hyd:euler-local}
 Let $s\ge4$ be an integer, and suppose
 $\theta_0\in H^s(\T^3;\R^5)$ has range compactly contained in
 $\mathcal O$. There exist $T_0>0$ and a unique solution
 \[
 \theta\in C([0,T_0];H^s)\cap C^1([0,T_0];H^{s-1})
 \]
 of \eqref{hyd:euler} with initial value $\theta_0$, whose range remains
 compactly contained in $\mathcal O$. Smooth initial data yield a smooth
 solution. The existence time is independent of $c$.
\end{proposition}

\begin{proof}
 Fix a compact neighborhood $K\Subset\mathcal O$ of the range of the
 initial data. If $q_-\le\theta\cdot\Psi\le q_+$, then, with
 $G_D=\int_D\Psi\Psi^{\mathsf T}\dd k$,
 \[
 q_+^{-2}G_D\le M(\theta)\le q_-^{-2}G_D.
 \]
 Consequently,
 $E_s=\sum_{|\alpha|\le s}\int (\partial^\alpha\theta)^{\mathsf T}
 M(\theta)\partial^\alpha\theta\dd X$ is equivalent to
 $\|\theta\|_{H^s}^2$. The multivariate Leibniz formula, the embedding
 $H^s\hookrightarrow W^{1,\infty}$, and interpolation give
 \begin{equation}\label{hyd:euler-comm}
 \sum_{|\alpha|\le s}
 \|[\partial^\alpha,A_j(\theta)]\partial_jz\|_2
 \le C_K\bigl(\|\nabla\theta\|_\infty\|z\|_{H^s}
                   +\|\theta\|_{H^s}\|\nabla z\|_\infty\bigr).
 \end{equation}
 Symmetry of $MA_j=J_j$ permits integration by parts in the leading
 transport term, so smooth solutions satisfy
 $|E_s'|\le C_K\|\nabla\theta\|_\infty E_s$.

 To construct a solution, let $S_\delta$ be convolution with an even
 periodic mollifier and solve the Banach-space ordinary differential
 equation
 \[
 (\theta_\delta)_t
 =-\sum_jS_\delta\{A_j(\theta_\delta)
                         \partial_jS_\delta\theta_\delta\},
 \qquad \theta_\delta(0)=\theta_0.
 \]
 For fixed $\delta$, the right-hand side is locally Lipschitz on the
 positive parameter domain in $H^s$. For
 $z=\partial^\alpha\theta_\delta$, split the principal term as
 \[
 \langle Mz,S_\delta(A_j\partial_jS_\delta z)\rangle
 =\langle MS_\delta z,A_j\partial_jS_\delta z\rangle
  +\langle[S_\delta,M]z,A_j\partial_jS_\delta z\rangle.
 \]
 Use symmetry for the first term and
 $\|[S_\delta,M]z\|_2\le C\delta\|\nabla M\|_\infty\|z\|_2$
 together with
 $\|\partial S_\delta z\|_2\le C\delta^{-1}\|z\|_2$ for the second.
 The same energy estimate therefore holds for the mollified equation,
 uniformly in $\delta$. The equation also gives
 $\|\theta_{\delta,t}\|_\infty
 \le C_K\|\nabla\theta_\delta\|_\infty$.
 Choose a short time on which both the energy and the positive distance
 to the boundary of $K$ are controlled. This gives a common existence
 interval.

 On this interval, $\theta_\delta$ is bounded in $L^\infty H^s$ and its
 time derivative in $L^\infty H^{s-1}$. Compact embedding and
 equicontinuity yield a subsequence converging strongly in
 $C_tH^{s-1}$, and the limit satisfies \eqref{hyd:euler}.
 Finally, mollify the differentiated equation for the limit once more.
 Integration by parts in the convolution kernel gives
 \[
 \|[S_\delta,A]\partial z\|_2
 \le C\|\nabla A\|_\infty\|z\|_2.
 \]
 Together with \eqref{hyd:euler-comm}, this gives a uniform integrable
 bound for the time derivative of the mollified energy. Letting
 $\delta\to0$ proves continuity of $E_s(t)$; weak $H^s$ continuity and
 continuity of the positive energy norm give strong $C_tH^s$ continuity.
 For the difference of two solutions, use the first solution's $M$ in
 an $L^2$ energy. Local Lipschitz continuity of the coefficients and the
 first derivatives of the other solution control the right-hand side;
 Gronwall's inequality proves uniqueness. Higher-order energies preserve
 smoothness successively.
\end{proof}

From now on, fix a smooth positive Euler solution $\bar\theta$ on a finite
interval $[0,T]$, with range compactly contained in $\mathcal O$, and write
$\bar N=N_{\bar\theta}$. This may be any interval of smooth positive
continuation of the solution in Proposition~\ref{hyd:euler-local}.
Constants may depend on its positivity margin and finitely many
derivative bounds. More precisely, fix a compact convex parameter
neighborhood $\Theta\Subset\mathcal O$ whose interior contains
$\bar\theta([0,T]\times\T^3)$, and set
\[
 K_5=1+\max_{j+|\alpha|\le5}
       \|\partial_t^j\partial_X^\alpha\bar\theta\|_{L^\infty([0,T]\times\T^3)}.
\]
The fixed data in the uniform estimates below are $R,T,\Theta,K_5$;
the temporary two-sided ratio bounds will also be determined by these
data. We use $C_K$ for a constant depending only on them, independent of
$\Lambda,\eta,c$. Dependence on an auxiliary parameter already chosen
will be indicated explicitly in the subscript. Higher derivatives of
the smooth initial data are used only for classical regularity at each
fixed $c$; they do not enter the uniform threshold $c_*$ or the constants
in the following three convergence rates.

\begin{theorem}[Strong-collision limit on a prescribed positive Euler interval]\label{hyd:main}
 Suppose $\theta_0\in C^\infty(\T^3;\R^5)$ has range compactly contained
 in $\mathcal O$, $\bar\theta(0)=\theta_0$, and $\bar\theta$ is the smooth
 positive Euler solution on $[0,T]$ specified above. There exist
 $c_*,C_T,m_T,M_T>0$ such that, for every $c\ge c_*$, the initial-value
 problem
 \begin{equation}\label{hyd:kinetic}
 \mathcal T f_c=c\mathcal C(f_c),\qquad f_c(0,X,k)=N_{\theta_0(X)}(k)
 \end{equation}
 has a unique positive solution on $[0,T]$. For every finite $m$ it
 belongs to $C_tC_X^mC_k(D)$; for $m\ge1$ it also belongs to
 $C_t^1C_X^{m-1}C_k(D)$, and $m_T\le f_c\le M_T$.
 Its five moments lie in $U(\mathcal O)$. Defining
 $N_c=N_{\theta_c}$ by these moments, we have
 \begin{align}
 \|f_c-\bar N\|_{L^\infty_tH_X^3L_k^2}
       +\|\theta_c-\bar\theta\|_{L^\infty_tH_X^3}
      &\le C_Tc^{-3/4},\label{hyd:main-H3}\\
 \|\theta_c-\bar\theta\|_{L^\infty_tH_X^2}
      &\le C_Tc^{-1},\label{hyd:main-H2}\\
 \|f_c-N_c\|_{L^2_tH_X^3H_\nu}&\le C_Tc^{-1}.
 \label{hyd:main-weighted}
 \end{align}
 Here $H_\nu=L^2(D,\nu_*\dd k)$, where $\nu_*$ is the collision
 frequency of any fixed positive Rayleigh--Jeans distribution.
 All constants are independent of $c$.
\end{theorem}

We distinguish five objects used in the proof. The prescribed Euler
solution determines $\bar N$, while the five moments of the evolving
solution $f_c$ determine $N_c$. The microscopic deviation is further
split into a corrector and a remainder:
\[
 \begin{aligned}
 \bar N&=N_{\bar\theta},
       &N_c&=N_{\theta_c},\qquad
          U(\theta_c)=\int_D\Psi f_c\,\dd k,\\
 f_c-N_c&=G_c+e_c,
       &\int_D\Psi G_c\,\dd k
          &=\int_D\Psi e_c\,\dd k=0.
 \end{aligned}
\]
The corrector $G_c$ will be constructed from $\bar N$ in
Subsection~\ref{hyd:cell-section}; $e_c$ is the distribution error
controlled by the joint energy. The difference between the two
macroscopic states is measured by $\theta_c-\bar\theta$.

A spatially homogeneous equilibrium gives the simplest instance of this
decomposition. If $\theta_0$ is constant, then
$\bar\theta(t,X)=\theta_0$, and $f_c(t,X,k)=N_{\theta_0}(k)$ is an exact
solution for every $c$: both transport and collision terms vanish.
The forcing in the corrector equation below also vanishes, its small
solution is $G_c=0$, and hence $e_c=0$.

The corners explain the choice of convergence topology. The collision
term vanishes at every corner $\xi$, so
\[
 f_c(t,X,\xi)=N_{\theta_0(X-2\xi t)}(\xi).
\]
This free-transport trace generally differs from $\bar N(t,X,\xi)$.
The momentum $L^2$ and weighted norms in the theorem allow this corner
memory; the corner comparison will provide positive two-sided bounds
in its neighborhood.

\subsection{Local construction and continuation of the transport equation}
\label{hyd:local-section}

\begin{lemma}[Amplitude continuation criterion]\label{hyd:local}
 Let $B_m=C_X^mC_k(D)$. For every fixed $c>0$ and strictly positive
 $f_0\in B_m$, equation~\eqref{hyd:kinetic} has a unique maximal positive
 mild solution. If its maximal existence time $T_{\max}$ is finite, then
 \[
 \sup_{t<T_{\max}}\|f(t)\|_\infty=\infty.
 \]
 For smooth spatial initial data, every finite $B_m$ norm remains finite
 on each finite interval of existence on which the amplitude is bounded.
\end{lemma}

\begin{proof}
 Lemma~\ref{geom:basic} gives local Lipschitz continuity of
 $\mathcal C:B_0\to B_0$ and
 $\|\mathcal C(f)\|_\infty\le C_R\|f\|_\infty^3$.
 The operators $S(t)f(X,k)=f(X-2kt,k)$ form a strongly continuous
 isometric group on $B_0$. The contraction mapping theorem applied to
 \[
 f(t)=S(t)f_0+c\int_0^tS(t-s)\mathcal C(f(s))\,\dd s
 \]
 on a bounded ball of continuous paths gives a unique local solution,
 with existence time depending only on $c$ and the amplitude bound of
 the ball. Reapplying the same local construction near a finite maximal
 endpoint extends the solution, and uniqueness identifies the extension
 with the original solution. This proves the amplitude criterion.

 Positivity follows from the decomposition of the original polynomial
 \begin{align*}
 \mathcal C(f)_0&=\mathcal G(f)_0+\mathcal B(f)_0 f_0,\\
 \mathcal G(f)_0&=\int f_1f_2f_3\,\dd\Gamma_{k_0}\ge0,\\
 \mathcal B(f)_0&=\int(f_2f_3-f_1f_2-f_1f_3)\,\dd\Gamma_{k_0}.
 \end{align*}
 Along transport characteristics, the equivalent integral equation
 obtained with the integrating factor for $c\mathcal B(f)$ is also a
 contraction on the closed nonnegative ball. Its solution therefore
 agrees with the one above. On an interval where the amplitude is at
 most $M$,
 $f(t)\ge(\min f_0)\exp(-cC_RM^2t)>0$.

 The group $S(t)$ commutes with $\partial_X$. In the full Leibniz
 expansion of each product of three factors, an $m$th-order derivative
 can fall entirely on only one factor. Hence
 \[
 \|f(t)\|_{B_m}\le\|f_0\|_{B_m}
   +cC_RM^2\int_0^t\|f(s)\|_{B_m}\,\dd s
   +c\int_0^tP_m(\|f(s)\|_{B_{m-1}})\,\dd s,
 \]
 where $P_m$ is a finite polynomial. Induction and Gronwall's inequality
 finish the proof. These higher-order norms at fixed $c$ justify
 differentiation; uniform energy bounds are proved separately below.
\end{proof}

\subsection{The nonlinear corrector in two norms}
\label{hyd:cell-section}

The local equilibrium satisfies $\mathcal C(\bar N)=0$, but its transport
residual $\mathcal T\bar N$ generally does not vanish. We construct a
corrector $G_c$ with zero five moments so that
$c\mathcal C(\bar N+G_c)$ cancels the leading part of this residual.
Degeneracy of the frequency at the corners calls for two distinct
controls: a small amplitude throughout momentum space and a stronger
$c^{-1}$ bound after multiplication by the frequency. Both norms will
be retained below.

We first construct and differentiate the equation in the Banach spaces
from the preceding section,
\[
 \mathcal X=L^\infty(D),\qquad
 \mathcal Y=\{h:\nu_*h\in L^\infty(D)\},
 \qquad\|h\|_{\mathcal Y}=\|\nu_*h\|_\infty,
\]
with functions identified almost everywhere. Write $X_0=C(D)$ for the
closed subspace of continuous functions. We will subsequently show that
continuous inputs give outputs and parameter derivatives in $X_0$.
For a general $\mathcal Y$ input, the diagonal term $\nu_Nh$ is only
known to be bounded; continuous representatives will be recovered after
the parameter estimates. For $s\ge1$, set
\[
 \|h\|_{\mathcal B_s}=\|h\|_\infty+s\|\nu_*h\|_\infty.
\]
For fixed $s$, the spaces $\mathcal B_s$ and $\mathcal X$ have the same
underlying set and equivalent norms. Its continuous subspace is again
denoted by $X_0$, equipped with the same $\mathcal B_s$ norm. Write
\begin{equation}\label{hyd:remainder}
 \widetilde{\mathcal C}_N(q)=N^{-1}\mathcal C(N(1+q))
             =-L_Nq+\mathcal R_N(q).
\end{equation}
The remainder $\mathcal R_N$ consists exactly of the quadratic and cubic
terms. Conservation of the five moments gives $P_N\mathcal R_N(q)=0$.

\begin{lemma}[The nonlinear cell problem and parameter derivatives]\label{hyd:cell}
 There is $\delta_0>0$, depending on $\Theta$, such that for real
 $F\in Q_NX_0$ with $\|F\|_\infty\le\delta_0$, the equation
 \begin{equation}\label{hyd:cell-small}
 q-s\widetilde{\mathcal C}_N(q)=F,\qquad P_Nq=0,
 \end{equation}
 has a unique real solution in a fixed small $\mathcal B_s$ ball, and
 $\|q\|_{\mathcal B_s}\le C_\Theta\|F\|_\infty$.
 Suppose $N=N(y)$ and $F=F(y)$ depend smoothly on a finite-dimensional
 parameter $y$, the specified finitely many derivatives of $N$ are
 bounded, and
 $\max_{|\alpha|\le j}\|\partial_y^\alpha F\|_\infty\le K_j\delta$.
 Then
 \begin{equation}\label{hyd:cell-jets}
 \max_{|\alpha|\le j}\bigl(
 \|\partial_y^\alpha q\|_\infty
       +s\|\nu_*\partial_y^\alpha q\|_\infty\bigr)
 \le C_j\delta.
 \end{equation}
 The constants are uniform for $s\ge1$.
\end{lemma}

\begin{proof}
 By Proposition~\ref{geom:resolvent},
 $R_{N,s}=(I+sL_N)^{-1}Q_N:\mathcal X\to\mathcal B_s$ is uniformly
 bounded. For $\|q_i\|_\infty\le\kappa$, Lemma~\ref{geom:remainder} gives
 \[
 \|\mathcal R_N(q_1)-\mathcal R_N(q_2)\|_\infty
       \le C\kappa\|\nu_*(q_1-q_2)\|_\infty.
 \]
 Thus $q\mapsto R_{N,s}[F+s\mathcal R_N(q)]$ is a contraction on the
 ball of radius $2C\|F\|_\infty$. First ensure $C\kappa<1/2$, then
 choose $\delta_0$. Iteration from zero remains in the same range of
 $Q_N$. The original cubic operator and the resolvent preserve
 continuous functions, so iteration with continuous input remains in
 $X_0$. The supremum and essential supremum of a continuous function on
 the closed cube agree; hence $X_0\hookrightarrow\mathcal X$ is a closed
 isometric embedding, and the limit is continuous. Iteration also
 preserves real values. Taking the ball radius below $1/2$ gives
 $1+q\ge1/2$.

 Parameter differentiation requires a fixed projected space. Retain the
 reference state $n_*$ and write $P_*=P_{n_*}$, $Q_*=I-P_*$.
 Use the operator $U_N$ from Lemma~\ref{geom:frame}, which satisfies
 $U_NP_*=P_NU_N$; the operators $U_N,U_N^*$ and their specified parameter
 derivatives are uniformly bounded on $\mathcal X,\mathcal Y$. Set
 $z=U_N^*q$, $\widetilde L_N=U_N^*L_NU_N$, and define
 \[
 \widetilde R_{N,s}=U_N^*R_{N,s}U_N
                  =(I+s\widetilde L_N)^{-1}Q_*,\qquad
 \widetilde{\mathcal R}_N(z)=U_N^*\mathcal R_N(U_Nz).
 \]
 On the fixed space $Q_*\mathcal X$, the resolvent satisfies
 \[
 \partial_\theta(I+s\widetilde L_N)^{-1}
  =-s(I+s\widetilde L_N)^{-1}(\partial_\theta\widetilde L_N)
                       (I+s\widetilde L_N)^{-1}.
 \]
 Reading the right-hand side from right to left gives an $s^{-1}$ gain
 from $\mathcal X$ to $\mathcal Y$, a bounded parameter derivative from
 $\mathcal Y$ to $\mathcal X$, and a uniform bound from $\mathcal X$ to
 $\mathcal B_s$. Every higher parameter derivative is controlled by the
 same pairing of factors $s$ and $s^{-1}$.

 In the nonlinear equation on the fixed space, the operator acting on
 the highest $y$ derivative is
 $I-s\widetilde R_{N,s}D\widetilde{\mathcal R}_N(z)$.
 Its inverse is uniformly controlled by the contraction constant above.
 Every remaining term in the finite-order chain rule contains parameter
 derivatives or lower state derivatives. Lemma~\ref{geom:remainder}
 permits exactly one input of each multilinear term to be estimated in
 $\mathcal Y$, with the others in $\mathcal X$. Each coefficient $s$ is
 therefore canceled by the $s^{-1}$ gain of a lower derivative. Induction
 gives \eqref{hyd:cell-jets}. Derivatives of $U_N^*$ are already included
 in this fixed-space equation; no separate membership of $F_y$ in the
 range of the original projection $Q_N$ is required.

 Finally, recover continuous representatives of the parameter
 derivatives. For fixed $s$, the argument above gives the required
 smoothness of $q(y)$ in $\mathcal X$. Its first difference quotients
 are continuous functions and converge in the $\mathcal X$ norm. The
 closed embedding gives a unique continuous representative of the
 limit, and the convergence is in fact uniform. Induction on higher
 difference quotients shows that all the stated derivatives belong to
 $X_0$ and depend continuously on the parameters. This recovery uses
 closedness only at each fixed $s$; the estimates uniform in $s$ have
 already been proved in $\mathcal X,\mathcal Y$.
\end{proof}

Apply this result with $F=-\mathcal T\log\bar N$. The Euler equation gives
\[
 \int_D\bar N\Psi F\,\dd k=-\int_D\Psi\mathcal T\bar N\,\dd k=0.
\]
The forcing $F$ may be large on the prescribed interval. Choose a fixed
$\Lambda\ge1$ and apply Lemma~\ref{hyd:cell} with $s=c/\Lambda$ and
$F/\Lambda$ to obtain
\begin{equation}\label{hyd:cell-Lambda}
 \Lambda q_c-c\widetilde{\mathcal C}_{\bar N}(q_c)
      =-\mathcal T\log\bar N,
 \qquad P_{\bar N}q_c=0,
 \qquad G_c=\bar Nq_c.
\end{equation}
Thus $\Lambda$ reduces the forcing in the corrector equation while the
positive parameter range of the initial data is unchanged. From now on,
$\Lambda$ depends on the prescribed interval, but not on $c$. For
$c\ge\Lambda\ge\Lambda_0$, all required time and space derivatives of
total order $|\alpha|\le4$ satisfy
\begin{equation}\label{hyd:G-two-norms}
 \|\partial^\alpha G_c\|_\infty\le C_K\Lambda^{-1},
 \qquad c\|\nu_*\partial^\alpha G_c\|_\infty\le C_K.
\end{equation}
The forcing $F=-\mathcal T\log\bar N$ uses one more derivative than
$\bar\theta$. Thus its total space-time derivatives through order four,
and the corresponding derivatives of $G_c$, are precisely controlled by
$K_5$. Taking the small parameter in the lemma to be $C_K/\Lambda$ and
$s=c/\Lambda$ gives both the $\Lambda^{-1}$ bound in the uniform norm
and the $c^{-1}$ bound in the weighted norm. The nonlinear products in
the parameter derivatives use estimates with one $\mathcal Y$ input.
The corner-frequency estimate in Lemma~\ref{geom:frequency} gives
$\nu_*^{-3/2}\in L^1(D)$. Since
$\min(A,B)^2\le A^{1/2}B^{3/2}$,
\begin{equation}\label{hyd:G-integral}
 \|\partial^\alpha G_c\|_2\le C_K\Lambda^{-1/4}c^{-3/4},
 \qquad
 \|\partial^\alpha G_c\|_{H_\nu}
       +\|\partial^\alpha G_c\|_1\le C_Kc^{-1}.
\end{equation}
The exponent $3/4$ comes from the critical integrability of the corner
frequency: squaring and integrating the pointwise bound
$|\partial^\alpha G_c|\le C_K\min\{\Lambda^{-1},(c\nu_*)^{-1}\}$
uses exactly $\nu_*^{-3/2}\in L^1$.
Let $f^{\mathrm{app}}=\bar N+G_c$. Its residual and initial error are
\begin{align}
 \mathcal E^{\mathrm{app}}
  &=\mathcal T f^{\mathrm{app}}-c\mathcal C(f^{\mathrm{app}})
    =\bar N\{\mathcal Tq_c+(\mathcal T\log\bar N)q_c-\Lambda q_c\},
       \label{hyd:residual}\\
 f_c(0)-f^{\mathrm{app}}(0)&=-G_c(0).
       \label{hyd:initial-layer}
\end{align}
In particular,
\begin{align}
 \|G_c\|_{L^\infty_tH_X^3L_k^2}
       &\le C_K\Lambda^{-1/4}c^{-3/4},\notag\\
 \|\mathcal E^{\mathrm{app}}\|_{L^\infty_tH_X^3L_k^2}
       &\le C_K\Lambda^{3/4}c^{-3/4},
 &\|\mathcal E^{\mathrm{app}}/\bar N\|_\infty&\le C_K.
       \label{hyd:residual-bounds}
\end{align}
The projection condition gives $\int\Psi G_c=0$, whereas
\begin{equation}\label{hyd:residual-moments}
 \int_D\Psi\mathcal E^{\mathrm{app}}\,\dd k
   =\sum_j\partial_j\int_Dv_j\Psi G_c\,\dd k.
\end{equation}
This vector is generally nonzero; the macroscopic equation below retains it.

\subsection{Orthogonal decomposition and differentiated equations through order three}
\label{hyd:coframe-section}

Work for now on an interval where the original solution exists. Use
$U_c=\int\Psi f_c$ to define $N=N_{\theta_c}$, and abbreviate
$G=G_c$, $f=f_c$, and $\theta=\theta_c$.
Throughout this subsection and the energy proof that follows, $N$
denotes the actual moment-matched equilibrium; the prescribed Euler
distribution remains $\bar N$. Set
\[
 e=f-N-G,\qquad a=\theta-\bar\theta,\qquad \kappa_N=N\Psi,
 \qquad \eta_\alpha=\frac{\partial_X^\alpha e}{N},
 \quad a_\alpha=\partial_X^\alpha a,
 \quad H_\alpha=\eta_\alpha-\kappa_N\cdot a_\alpha.
\]
These adapted derivative variables are obtained by first differentiating
the distribution error $e$ in space and then dividing by $N$.
In particular,
\[
 \eta_j=\partial_j(e/N)+(\partial_j\log N)(e/N).
\]
This choice preserves exact five-moment orthogonality at each order;
variation of the denominator appears in the connection terms below.
The macroscopic and microscopic components thus use the same transport
velocity, and the highest-order cross terms cancel in the joint energy.
For every $0\le|\alpha|\le3$,
\begin{equation}\label{hyd:orthogonality}
 \int_D\kappa_N\eta_\alpha\,\dd k
       =\int_D\Psi\partial^\alpha e\,\dd k=0.
\end{equation}
Define
\begin{align}
 E_l&=\frac12\sum_{|\alpha|=l}\int_{\T^3}
       (\|\eta_\alpha\|_2^2+a_\alpha^{\mathsf T}Ma_\alpha)\,\dd X
     =\frac12\sum_{|\alpha|=l}\|H_\alpha\|_{L^2_{X,k}}^2,
       \label{hyd:E-l}\\
 D_l&=\sum_{|\alpha|=l}\int_{\T^3}
                 \mathfrak a_N(\eta_\alpha,\eta_\alpha)\,\dd X,
 &E&=\sum_{l=0}^3\omega_lE_l,
 &D_e&=\sum_{l=0}^3\omega_lD_l.\label{hyd:D-l}
\end{align}
The fixed decreasing weights
$1=\omega_0>\omega_1>\omega_2>\omega_3>0$ will be chosen in the energy
estimate. In the parameter neighborhood where $E\le e_*$,
Proposition~\ref{geom:coercivity} and \eqref{hyd:orthogonality} give
\begin{equation}\label{hyd:energy-equivalence}
 E\asymp\|e\|_{H_X^3L_k^2}^2+\|a\|_{H_X^3}^2,
 \qquad \|e\|_{H_X^lH_\nu}^2\le C\sum_{p=0}^lD_p.
\end{equation}

We now derive the differentiated identities, including order zero. Set
\[
 B=\mathcal T\log N,\qquad
 C_\alpha=\frac{\partial^\alpha N}{N}
                      +\kappa_N\cdot\partial^\alpha\theta,
 \qquad s_\alpha=\kappa_N\cdot\partial^\alpha\theta.
\]
Here $\partial_j=\partial/\partial X_j$. In $C_{ij}$ and $s_{ij}$ the
indices indicate the successive spatial derivatives
$\partial_i\partial_j$, and repeated indices are allowed. Sums in the
energy are still indexed by multiindices $\alpha$. Direct differentiation
gives
\begin{align}
 C_0&=2,& C_i&=0,& C_{ij}&=2s_is_j,\notag\\
 C_{ijk}&=2(s_is_{jk}+s_js_{ik}+s_ks_{ij})-6s_is_js_k.
 \label{hyd:B-jets}
\end{align}
Consequently,
\begin{equation}\label{hyd:H-identity}
 H_\alpha=\frac{\partial^\alpha(f-G)}N
                 +\kappa_N\cdot\partial^\alpha\bar\theta-C_\alpha.
\end{equation}
Set
\begin{align}
 S_\alpha&=\kappa_N\cdot\mathcal T\partial^\alpha\bar\theta
       +2B\kappa_N\cdot\partial^\alpha\bar\theta
       -(\mathcal T+B)C_\alpha,\notag\\
 \mathcal R_\alpha^{\mathrm{src}}&=S_\alpha+\frac{\partial^\alpha\mathcal T\bar N}{N}.
 \label{hyd:R-definition}
\end{align}
Differentiating \eqref{hyd:H-identity} and using \eqref{hyd:residual}
yields
\begin{equation}\label{hyd:H-equation}
 \mathcal TH_\alpha
 =\frac cN\partial^\alpha[\mathcal C(f)-\mathcal C(f^{\mathrm{app}})]
      -BH_\alpha+\mathcal R_\alpha^{\mathrm{src}}-
       \frac{\partial^\alpha\mathcal E^{\mathrm{app}}}{N}.
\end{equation}
Its collision term is orthogonal to $\kappa_N$. In particular, $C_0=2$
gives
\begin{equation}\label{hyd:R-zero}
 \mathcal R_0^{\mathrm{src}}=(1-\rho^2)\kappa_N\cdot\mathcal T\bar\theta
          +2B(\kappa_N\cdot\bar\theta-1),\qquad \rho=\bar N/N.
\end{equation}

The actual five-moment equation is
\begin{equation}\label{hyd:actual-macro}
 \theta_t=-\sum_jA_j(\theta)\partial_j\theta
                 +M^{-1}\sum_j\partial_jZ_j,
 \qquad Z_j=\int_Dv_j\Psi(e+G)\,\dd k.
\end{equation}
It is obtained by taking moments of the conservation law for the same
distribution; in particular, it contains no collision coefficient $c$.
Split the flux into the distribution-error and known-corrector parts:
\[
 Z_j^e=\int_Dv_j\Psi e\,\dd k,\qquad
 Z_j^G=\int_Dv_j\Psi G\,\dd k,\qquad Z_j=Z_j^e+Z_j^G.
\]
Spatial derivatives commute with this moment integral, so
\[
 \|Z\|_{H_X^3}\le
 C\|e\|_{H_X^3L_k^2}+C\|G\|_{H_X^3L_k^1}
 \le C_K(\sqrt E+c^{-1}).
\]
Subtract the Euler equation from \eqref{hyd:actual-macro} and use the
product rule in $H^2_X$ to obtain the first estimate below. Three spatial
derivatives suffice for $\operatorname{div}Z$; third-order control of
$\theta_t$ is not required. By \eqref{hyd:G-integral},
\eqref{hyd:energy-equivalence}, and $H^3_X\hookrightarrow C^1_X$,
\begin{align}
 \|\theta_t-\bar\theta_t\|_{H_X^2}
       +\|B-\bar B\|_{H_X^2L_k^\infty}
       &\le C_K(\sqrt E+c^{-1}),\notag\\
 \|B\|_\infty&\le C_K,
 &\sum_{|\alpha|\le3}\|\mathcal R_\alpha^{\mathrm{src}}\|_{L_X^2L_k^\infty}
       &\le C_K(\sqrt E+c^{-1}),
 \label{hyd:connection-bounds}
\end{align}
where $\bar B=\mathcal T\log\bar N$. To explain the highest-order case
of the last estimate, note that for $1\le|\beta|\le2$,
\[
 \mathcal Ts_\beta
 =B s_\beta+\kappa_N\cdot
       \left(\partial^\beta\theta_t+
             \sum_jv_j\partial_j\partial^\beta\theta\right).
\]
Substituting this into the third-order formula in \eqref{hyd:B-jets},
the highest time-derivative products are $\theta_{tij}\theta_i$ and
$\theta_{ti}\theta_{jk}$; the highest spatial terms are
$\theta_{ijk}\theta_i$ and products of two second derivatives.
Estimate $\theta_{tij}\theta_i$ and $\theta_{ijk}\theta_i$ in
$L^2\times L^\infty$, and the remaining products in $L^6\times L^3$.
All these norms are controlled by $\theta_t\in H^2$ and $\theta\in H^3$.
At $\theta=\bar\theta$,
$S_\alpha=-\partial^\alpha\mathcal T\bar N/\bar N$; hence every term
after subtraction contains either a parameter error or a microscopic
moment from \eqref{hyd:actual-macro}. Equation~\eqref{hyd:R-zero}
verifies the zeroth-order endpoint of the same argument.

\subsection{Third-order energy under two-sided corner-layer bounds}
\label{hyd:energy-section}

Let $E_\eta$ be the union of the $\eta$-neighborhoods of the eight corners,
with distance measured in the $\ell^\infty$ norm. On a temporary interval
of existence, assume
\begin{equation}\label{hyd:bootstrap-assumptions}
 0<d_-\le f/N\le d_+<\infty,
 \qquad \sup_{k\notin E_\eta}|f/N-1|\le\kappa_b,
 \qquad E\le e_*.
\end{equation}
The constants $d_\pm$ are fixed, and deviations within the corner layers
need not be small. Write $v_\eta=\eta^2[1+\log(R/\eta)]^2$. If
\begin{equation}\label{hyd:bulk-G-condition}
 c v_\eta\ge C_K/\kappa_b,
\end{equation}
then \eqref{hyd:G-two-norms} and frequency comparison give
$|e|\le C\kappa_b$ on $D\setminus E_\eta$, while $|e|\le K_0$ everywhere,
where $K_0$ depends only on the fixed positive two-sided bounds. This
converts the assumptions on the distribution $f$ into assumptions on
the error $e$ used in the differentiated estimates.

\begin{proposition}[Joint energy estimate]\label{hyd:energy}
 On the interval above, there are fixed decreasing order weights and
 constants $\gamma,C_K>0$ such that, for sufficiently small
 $\kappa_b,\eta,e_*$ and sufficiently large $\Lambda$, provided
 \eqref{hyd:bulk-G-condition} holds,
 \begin{equation}\label{hyd:energy-differential}
 E'+\gamma cD_e\le C_KE+C_K\Lambda^{3/2}c^{-3/2}.
 \end{equation}
 The smallness threshold can be written as
 \begin{equation}\label{hyd:energy-smallness}
 C_K\bigl(\kappa_b^{1/3}+\omega_\eta^{1/3}
                      +\Lambda^{-1}+\sqrt{e_*}\bigr)\le\gamma_0,
 \end{equation}
 where $\omega_\eta\to0$ is the corner-restricted two-leg operator norm
 in Lemma~\ref{geom:moser}. Constants may depend on $d_\pm$, the positive
 parameter domain, and finitely many derivatives of the prescribed
 Euler solution, but not on $c,\eta,\Lambda$ or any uncontrolled
 higher-order norm of the actual solution.
\end{proposition}

The estimates play distinct roles. The highest collision derivatives
produce $D_l$; parameter variations and the corrector produce terms
controlled by lower-order dissipation and the known residual. The
actual five-moment equation handles the transport connection terms.
Fixed decreasing order weights then absorb the lower-triangular
dissipation terms successively. The four small quantities in
\eqref{hyd:energy-smallness} correspond respectively to the error outside
the corner layers, two-leg interactions in the corner layers, the
corrector amplitude, and the unknown macroscopic parameter error.

\begin{proof}
 The proof separates collision differentiation, parameter variation,
 and transport. Let $F=N+G$ and write
 \begin{equation}\label{hyd:collision-split}
 \mathcal C(f)-\mathcal C(f^{\mathrm{app}})
 =\underbrace{\mathcal C(F+e)-\mathcal C(F)}_{T(e;N,G)}
  +\underbrace{\mathcal C(N+G)-\mathcal C(\bar N+G)}_{V(a;G)}.
 \end{equation}

 First consider $T$. The distributions $F+se$, $0\le s\le1$, have the
 same five moments and the same moment-matched equilibrium $N$. They
 are positive, have uniform two-sided bounds, and are close to $N$
 outside the corner layers. Proposition~\ref{geom:localized} and
 \eqref{hyd:orthogonality} give
 \begin{equation}\label{hyd:collision-zero}
 \left\langle\eta_0,\frac{T}{N}\right\rangle
 =\int_0^1\left\langle\eta_0,
      \frac{D\mathcal C(F+se)[N\eta_0]}N\right\rangle\dd s
 \le-\gamma_1D_0.
 \end{equation}
 This supplies the negative term in the zeroth-order energy directly.

 For positive order $|\alpha|=l$, the chain rule for the cubic polynomial
 gives
 \begin{equation}\label{hyd:highest-split}
 \partial^\alpha T
 =D\mathcal C(f)[\partial^\alpha e]
   +[D\mathcal C(f)-D\mathcal C(F)]\partial^\alpha F
   +\mathcal L_\alpha.
 \end{equation}
 The coefficients of the first term retain all values of $f$ in the
 corner layers, including those not close to equilibrium. Nevertheless,
 \eqref{hyd:orthogonality} and Proposition~\ref{geom:localized} yield
 \begin{equation}\label{hyd:highest-negative}
 \left\langle\eta_\alpha,
            \frac{D\mathcal C(f)[N\eta_\alpha]}N\right\rangle
       \le-\gamma_1\mathfrak a_N(\eta_\alpha,\eta_\alpha).
 \end{equation}

 We estimate all remaining terms by type. The four signs and leg-index
 sets in the collision polynomial are
 \[
 (+,123),\quad (+,023),\quad (-,013),\quad (-,012).
 \]
 For each such index set $I$, the full derivative formula is
 \begin{equation}\label{hyd:all-parents}
 \partial^\alpha\prod_{i\in I}f_i
 =\sum_{\sum_{i\in I}\beta_i=\alpha}
          \frac{\alpha!}{\prod_{i\in I}\beta_i!}
             \prod_{i\in I}\partial^{\beta_i}f_i.
 \end{equation}
 Substitute $f=N+G+e$ and $F=N+G$ separately and subtract. Every remaining
 term contains at least one factor $e$. Equation~\eqref{hyd:highest-split}
 separates the terms in which a single factor receives the entire
 multiindex $\alpha$; $\mathcal L_\alpha$ is the difference of the
 remaining multiindex terms.

 Fix a positive $n_*$ and use the common four-leg measure
 $\dd\mu_*=\prod_i n_*(k_i)\dd\Xi_D$. Each one-leg marginal is
 $n_*^2\nu_*\dd k$. First compare the varying $N$ weights with this
 fixed measure. Lemma~\ref{geom:moser} applies directly to the
 distribution error $e$: momentum indicators commute with every
 $X$ derivative, and \eqref{hyd:bulk-G-condition} gives the required
 amplitude bound outside the corner layers. Write
 $R_l=\|e\|_{H_X^lH_\nu}$. For total order $l=2,3$, terms in which
 at least two positive-order derivatives act on distinct $e$ factors
 satisfy
 \begin{equation}\label{hyd:localized-product}
 \|\text{these products}\|_{L_X^2L^2(\mu_*)}
 \le C_{K_0}(\kappa_b^{1/3}+\omega_\eta^{1/3})R_l.
 \end{equation}
 Any zeroth-order factor in a product of three factors is estimated
 only by $K_0$. If the total derivative order on $e$ is at most $l-1$,
 the ordinary product estimate on the same measure gives
 $C_{K_0}R_{l-1}$.

 To specify how these estimates apply here, divide the remainder into
 three types. The first contains no $G$, and all $l$ derivatives fall
 on $e$. Apart from the leading term already separated,
 \eqref{hyd:localized-product} applies exactly. The second contains no
 $G$, but at least one derivative falls on $N$. Terms containing only
 derivatives of the prescribed reference parameters contribute
 $C_KR_{l-1}$. Subtract the remaining smooth parameter coefficients
 along $\bar\theta+sa$; their differences contain at least one derivative
 of $a$. Through total order three,
 \begin{align*}
 \|a_3 e_0\|_{L_X^2H_\nu}
      &\le\|a_3\|_{L_X^2}\|e\|_{L_X^\infty H_\nu}
       \le C\|a\|_{H_X^3}R_3,\\
 \|a_2 e_1\|_{L_X^2H_\nu}
      &\le\|a_2\|_{L_X^6}\|e_1\|_{L_X^3H_\nu}
       \le C\|a\|_{H_X^3}R_3,\\
 \|a_1 e_2\|_{L_X^2H_\nu}
      &\le C\|a\|_{H_X^3}R_3.
 \end{align*}
 An additional $e$ factor is handled by its amplitude bound or a product
 estimate on the same measure. All parameter coefficients are expanded
 according to total derivative order. These three cases and their
 lower-order versions therefore exhaust products of highest parameter
 derivatives with microscopic errors, giving $C_K\sqrt E R_3$.
 In particular, the second term in \eqref{hyd:highest-split} obeys the
 same estimate: indeed,
 \[
 \partial^\alpha N=-N\kappa_N\cdot\partial^\alpha\theta+NC_\alpha,
 \]
 and every coefficient of $D\mathcal C(f)-D\mathcal C(F)$ contains an
 $e$ factor. A prescribed highest derivative
 $\partial^\alpha\bar\theta$ times a lower-order $e$ contributes
 $C_KR_{l-1}$; an unknown highest derivative $\partial^\alpha a$
 contributes $C_K\sqrt E R_3$ by the estimates above.

 The third type contains one or more factors $G$. Every required
 derivative of $G$ is bounded in the uniform norm by $C_K/\Lambda$
 through \eqref{hyd:G-two-norms}. The remaining terms still contain $e$,
 so the ordinary bounded-amplitude product estimate on the same measure
 gives $C_K\Lambda^{-1}R_3$; terms containing a parameter error add
 $C_K\sqrt E R_3$. This classification retains every sign and
 multiplicity in \eqref{hyd:all-parents}, including products involving
 the output leg.

 More explicitly, among the remaining terms without $G$, let $r_N$ be
 the total derivative order on the parameter coefficients. If $r_N=0$,
 the only nonleading allocations of third-order microscopic derivatives
 are $2+1$ and $1+1+1$, and the only second-order allocation is $1+1$.
 If $r_N\ge1$, the microscopic derivative order is at most $l-1$.
 After subtracting the prescribed reference values of the parameter
 coefficients, the highest unknown allocations are precisely
 $a_3e_0,a_2e_1,a_1e_2$ and their lower-order or multifactor variants.
 This includes a zeroth-order $e$ factor times a highest parameter
 derivative; no additional $c^{-1}$ factor is required on that
 highest parameter derivative.

 The one-leg marginal norm of the test function $\eta_\alpha(k_0)$ is
 controlled by $D_l^{1/2}$. For any product $P$ above, after removing
 uniformly bounded weights depending only on the positive parameter
 domain, the estimate used is
 \[
 \left|\int_{\T^3}\int\eta_\alpha(k_0)P\,\dd\mu_*\dd X\right|
 \le \|\eta_\alpha(k_0)\|_{L_X^2L^2(\mu_*)}
                         \|P\|_{L_X^2L^2(\mu_*)}.
 \]
 The marginal measure in the first factor is exactly
 $n_*^2\nu_*\dd k_0$. Thus all three classes of product bounds can be
 paired with $D_l^{1/2}$, without integrating the four legs separately
 or assuming independence. Applying Young's inequality to the known
 lower-order terms gives
 \begin{align}
 \sum_{|\alpha|=l}\left\langle\eta_\alpha,
                          \frac{\partial^\alpha T}{N}\right\rangle
 &\le-\gamma_2D_l+C_l\sum_{p<l}D_p\notag\\
 &\quad+C_K(\kappa_b^{1/3}+\omega_\eta^{1/3}
                    +\Lambda^{-1}+\sqrt E)\sum_{p=0}^3D_p.
 \label{hyd:triangular-collision}
 \end{align}
 For $l=0$, the lower-order sum is empty and
 \eqref{hyd:collision-zero} applies directly. To choose the order
 weights, first fix these constants on the ball $\|a\|_{H^3}\le1$,
 retaining all unknown parameter errors in the $\sqrt E$ term.
 Set $\omega_0=1$ and choose successively
 $0<\omega_l<\omega_{l-1}/2$ so that
 \[
 \omega_l C_l\le\frac{\gamma_2}{8}
                  \min_{0\le p<l}\omega_p,\qquad l=1,2,3.
 \]
 Each $D_p$ receives at most three such lower-order contributions,
 which are absorbed by its negative term. Decrease the threshold in
 \eqref{hyd:energy-smallness} so that the common remainder is at most
 $\gamma_2\sum_p\omega_pD_p/4$. Finally, choose $e_*$ so that
 \eqref{hyd:energy-equivalence} implies $\|a\|_{H^3}\le1$ and the
 smallness condition holds. Thus the order weights are fixed before
 the energy threshold, and neither depends on time or $c$.
 The contribution of $T$ is at most $-\gamma cD_e$.

 For $V$, vary $\theta_s=\bar\theta+sa$ while keeping the distribution
 corrector $G$ fixed. Since $\mathcal C(N_{\theta_s})=0$, every term
 obtained after any fixed number of parameter derivatives still
 contains at least one $G$ factor. If that factor lies on the output
 leg, the collision integral gives $C\nu_NG$; if it lies on an internal
 leg, the bound
 $\int K_{\mathrm a,N}(k,p)/\nu_*(p)\dd p\le C$ from
 Lemma~\ref{geom:rows} gives the same estimate. In each term, estimate
 exactly one $G$ in $\mathcal Y$ and the others in the uniform norm.
 Since $c\|\nu_*\partial^\beta G\|_\infty\le C_K$, the multivariate
 chain rule yields
 \begin{equation}\label{hyd:macro-collision-variation}
 c\|\mathcal C(N+G)-\mathcal C(\bar N+G)\|_{H_X^3L_k^2}
       \le C_K\|a\|_{H_X^3}.
 \end{equation}
 The relation between parameter differentiation and normalization can
 be displayed explicitly. Let $N_s=N_{\theta_s}$, $q_s=G/N_s$, and
 $b_s=\partial_s\log N_s=-N_s(a\cdot\Psi)$. Write $D_N$ for the
 Fr\'echet derivative with respect to the positive distribution $N$
 while $q$ is held fixed, namely
 \[
 D_N\widetilde{\mathcal C}_N(q)[h]
  =-\frac hN\widetilde{\mathcal C}_N(q)
    +\frac1N D\mathcal C(N(1+q))[h(1+q)].
 \]
 Since $\partial_sN_s=N_sb_s$ and $\partial_sq_s=-b_sq_s$, the chain
 rule gives
 \[
 \frac{\partial_s\mathcal C(N_s+G)}{N_s}
 =b_s\widetilde{\mathcal C}_{N_s}(q_s)
  +D_N\widetilde{\mathcal C}_{N_s}(q_s)[N_sb_s]
  -D_q\widetilde{\mathcal C}_{N_s}(q_s)[b_sq_s].
 \]
 All three terms arise from differentiation with the same distribution
 corrector $G$ held fixed. Pairing
 \eqref{hyd:macro-collision-variation} with the unweighted energy
 contributes only $C_KE$.

 Finally, integrate \eqref{hyd:H-equation} over the periodic spatial
 domain. The collision term has zero five moments, so its pairing with
 $H_\alpha$ can be replaced by the pairing with $\eta_\alpha$. Thus
 \begin{align}
 E'={}&c\sum_{l,|\alpha|=l}\omega_l
  \left\langle\eta_\alpha,
     \frac{\partial^\alpha[\mathcal C(f)-\mathcal C(f^{\mathrm{app}})]}{N}
  \right\rangle\notag\\
 &-\sum_{l,|\alpha|=l}\omega_l\int B|H_\alpha|^2\,\dd X\dd k\notag\\
 &+\sum_{l,|\alpha|=l}\omega_l
  \left\langle H_\alpha,\mathcal R_\alpha^{\mathrm{src}}
               -\frac{\partial^\alpha\mathcal E^{\mathrm{app}}}{N}\right\rangle.
 \label{hyd:exact-energy}
 \end{align}
 By \eqref{hyd:connection-bounds} and the ordinary $L^2$ Young
 inequality, the last two lines are bounded by
 $C_KE+C_Kc^{-2}+C_K\Lambda^{3/2}c^{-3/2}$.
 Combining this with the collision estimates proves the result.
 The residual term $\Lambda q_c$ and the zeroth-order initial layer
 have both been retained throughout.
\end{proof}

The initial data $f(0)=N(0)=\bar N(0)$ give $a(0)=0$ and $e(0)=-G(0)$.
Proposition~\ref{hyd:energy} and \eqref{hyd:G-integral} therefore yield
\begin{equation}\label{hyd:enhanced-bound}
 \sup_{s\le t}E(s)+c\int_0^tD_e(s)\,\dd s
 \le C_K e^{C_KT}(\Lambda^{-1/2}+T\Lambda^{3/2})c^{-3/2},
 \qquad t\le T,
\end{equation}
whenever the hypotheses hold on $[0,t]$. The right-hand side uses the
prescribed $T$ and is uniform on every shorter interval; it contains no
negative power of that interval's length.

\subsection{Positive exchange in corner layers and amplitude bounds outside them}
\label{hyd:corner-section}

\begin{lemma}[Corner-layer comparison]\label{hyd:corner}
 Assume \eqref{hyd:bootstrap-assumptions} and
 \eqref{hyd:enhanced-bound} hold on $[0,t]$. For each corner $\xi$, set
 \begin{gather*}
 F_\xi(s,X)=N_{\theta_0(X-2\xi s)}(\xi),\qquad
 F_-\le F_\xi\le F_+,\\
 A_* =\min(1,N_-/(2F_+)),\qquad
 B_* =\max(1,2N_+/F_-),
 \end{gather*}
 where $N_-\le N_\theta\le N_+$ uniformly on the parameter domain in use.
 Then
 \begin{equation}\label{hyd:corner-barrier}
 A_*-r_c^{\mathrm{cor}}\le\frac{f(s,X,k)}{F_\xi(s,X)}\le B_*+r_c^{\mathrm{cor}},
 \quad k\in E_\eta(\xi),\ s\le t,
 \end{equation}
 and
 \begin{equation}\label{hyd:corner-remainder}
 r_c^{\mathrm{cor}}\le C_{K,\Lambda,T,d_\pm}
     [\eta+c\eta v_\eta+\psi(\eta)+c^{-1/4}v_\eta^{-1/4}],
 \qquad\psi(\eta)\longrightarrow0.
 \end{equation}
\end{lemma}

\begin{proof}
 Fix $k_0\in E_\eta(\xi)$ and partition the original integration
 variables into
 \begin{align*}
 \mathcal B&=\{k_1\in E_{4\eta}\}
       \cup\{k_2\text{ or }k_3\text{ lies in }E_\eta\text{ at another corner}\},\\
 \mathcal A_2&=\mathcal B^c\cap\{k_2\in E_\eta(\xi)\},\\
 \mathcal A_3&=\mathcal B^c\cap\{k_3\in E_\eta(\xi)\},\\
 \mathcal A_0&=\mathcal B^c\cap\{k_2,k_3\notin E_\eta\}.
 \end{align*}
 If both $k_2,k_3$ belong to $E_\eta(\xi)$, momentum conservation forces
 $k_1\in E_{3\eta}$, so $\mathcal A_2$ and $\mathcal A_3$ are disjoint.
 Lemma~\ref{geom:roles} gives
 $\int_{\mathcal B}\dd\Gamma_{k_0}\le C\eta\nu_*(k_0)$.
 On $\mathcal A_2$, the sum of the four cubic terms is
 \[
 f_1f_3(f_2-f_0)+f_0f_2(f_3-f_1),
 \]
 and the corresponding expression on $\mathcal A_3$ interchanges
 indices $2,3$. The first term defines an exchange generator with
 nonnegative rates:
 \begin{align*}
 (\mathcal J_f h)(k_0)
  ={}&\int_{\mathcal A_2}f_1f_3[h(k_2)-h(k_0)]\,\dd\Gamma_{k_0}\\
    &+\int_{\mathcal A_3}f_1f_2[h(k_3)-h(k_0)]\,\dd\Gamma_{k_0}.
 \end{align*}
 It exchanges values only at the same $X$ and within the same corner
 layer, and $\mathcal J_f1=0$.

 On $\mathcal A_0$, all three internal values are close to $N$.
 Therefore, for sufficiently small $\kappa_b$,
 \[
 a_0(f)=f_1f_3+f_1f_2-f_2f_3>0,
 \qquad \frac{f_1f_2f_3}{a_0(f)}\in[N_-/2,2N_+].
 \]
 Set $\nu_b=\int_{\mathcal A_0}a_0(f)\dd\Gamma_{k_0}$ and
 $\nu_b F_b=\int_{\mathcal A_0}f_1f_2f_3\dd\Gamma_{k_0}$.
 When $\nu_b=0$, choose any $F_b$ in the indicated interval.
 The original equation is then equivalent to
 \begin{equation}\label{hyd:positive-exchange}
 \mathcal T f=c\mathcal J_f f+c\nu_b(F_b-f)
                       +c(S_N+S_G+S_e+S_{\mathcal B}).
 \end{equation}
 Here $S_{\mathcal B}$ is the entire original polynomial integrated
 over the exceptional region. The other three source terms come from
 the decomposition
 $f_3-f_1=(N_3-N_1)+(G_3-G_1)+(e_3-e_1)$ on $\mathcal A_2$ and the
 corresponding expression on $\mathcal A_3$.

 With the positive two-sided bounds fixed,
 $c\|S_{\mathcal B}\|_{L^1_sL^\infty_{X,k}(E_\eta)}
 \le C_Tc\eta v_\eta$. On $\mathcal A_2$,
 $k_1-k_3=k_2-k_0=O(\eta)$, and the momentum Lipschitz constant of
 $N_\theta$ is uniform on the parameter domain. Thus $S_N$ obeys the
 same bound. Lemma~\ref{geom:rows} gives
 \[
 \sup_{k\in E_\eta,\,\theta\in\Theta}
    \int_D\frac{K_{\mathrm a,\theta}(k,p)}{\nu_*(p)}\,\dd p
       =:\psi(\eta)\longrightarrow0.
 \]
 Together with $c\|\nu_*G\|_\infty\le C_K$, this yields
 $c\|S_G\|_{L^1_sL^\infty_{X,k}(E_\eta)}\le C_T\psi(\eta)$.

 In $S_e$, the factors $e_1,e_3$ lie in $D\setminus E_\eta$.
 The square-row estimate of the same lemma gives
 \[
 \int_D K_{\mathrm a,\theta}(k,p)^2\,\dd p
       \le C\sqrt{\nu_*(k)}.
 \]
 Cauchy's inequality and
 $H_X^3(H_\nu)\hookrightarrow L_X^\infty(H_\nu)$ therefore imply
 \[
 \begin{split}
 \left|\int_{D\setminus E_\eta}K_{\mathrm a,\theta}(k,p)
                                      e(s,X,p)\,\dd p\right|
 &\le C\nu_*(k)^{1/4}v_\eta^{-1/2}
                                      \|e(s,X)\|_{H_\nu}\\
 &\le Cv_\eta^{-1/4}\|e(s,X)\|_{H_\nu},\qquad k\in E_\eta.
 \end{split}
 \]
 The first line uses $\nu_*(p)\ge c_0v_\eta$ on the input region and
 the square-row estimate; the second uses $\nu_*(k)\le C_0v_\eta$ on
 the output region. Consequently,
 \[
 \sup_{X,k\in E_\eta}|S_e(s,X,k)|
 \le Cv_\eta^{-1/4}\|e(s)\|_{H_X^3H_\nu}.
 \]
 Take this joint supremum before integrating in time. The bound
 $\int_0^t\|e(s)\|_{H_X^3H_\nu}^2\dd s\le C_{K,\Lambda,T}c^{-5/2}$
 in \eqref{hyd:enhanced-bound}, followed by Cauchy's inequality in time,
 gives
 \begin{equation}\label{hyd:corner-error-source}
 c\|S_e\|_{L^1_sL^\infty_{X,k}(E_\eta)}
       \le C_{K,\Lambda,T}c^{-1/4}v_\eta^{-1/4}.
 \end{equation}

 The collision row vanishes at a corner, so
 $f(s,X,\xi)=F_\xi(s,X)$. For nearby momenta,
 $\mathcal T\log F_\xi
 =2(k-\xi)\cdot\nabla_X\log F_\xi=O(\eta)$.
 Set $y=f/F_\xi$. Since exchange occurs at the same $X$ and within
 the same corner layer, \eqref{hyd:positive-exchange} becomes
 \[
 \mathcal Ty=c\mathcal J_fy+c\nu_b(F_b/F_\xi-y)
   +c(S_N+S_G+S_e+S_{\mathcal B})/F_\xi
   -y\mathcal T\log F_\xi.
 \]
 Fix the coefficients along this solution. Transport together with
 exchange preserves positivity and constants; after adding the
 nonnegative killing rate $c\nu_b$, the evolution is contractive in
 $L^\infty$. Explicitly, incorporate the loss rate into an integrating
 factor along characteristics and iterate the positive gain term to
 obtain a positive evolution. Comparison with constants shows that its
 norm is at most one. The relaxation target
 $F_b/F_\xi$ lies in $[A_*,B_*]$. Initially $y(0)=1+O(\eta)$, and the
 last term is $O(\eta)$ by the temporary amplitude bound. The preceding
 $L^1_sL^\infty$ estimates now give
 \eqref{hyd:corner-barrier}--\eqref{hyd:corner-remainder}.
\end{proof}

\begin{lemma}[Amplitude bound outside the corner layers]\label{hyd:bulk}
 On the same interval as in Lemma~\ref{hyd:corner}, set
 $h=(f-f^{\mathrm{app}})/\bar N$. If the actual parameter error is
 sufficiently small, then
 \begin{equation}\label{hyd:bulk-bound}
 \sup_{s\le t,\,X,\,k\notin E_\eta}|h(s,X,k)|
 \le C_{K,\Lambda,T}\{(cv_\eta)^{-1}
                           +c^{-3/4}v_\eta^{-3/4}\}.
 \end{equation}
\end{lemma}

\begin{proof}
 Apply Proposition~\ref{geom:localized} to the linear interpolation
 from $f^{\mathrm{app}}$ to $f$. Relative to $\bar N$, these distributions
 have fixed positive two-sided bounds, and their deviations outside
 the corner layers are controlled by
 $C\kappa_b+C_K/\Lambda+C\sqrt E$. The resulting equation is
 \[
 \mathcal Th+\bar B h+c\nu_{\mathrm{eff}}h
       =-cK_{\mathrm{eff}}h-\mathcal E^{\mathrm{app}}/\bar N,
 \qquad h(0)=-q_c(0),
 \]
 where $\nu_{\mathrm{eff}}\asymp\nu_*$ and all three signed
 off-diagonal kernels are dominated by $C K_{\mathrm a,\bar N}$.
 Outside the corner layers, when $cv_\eta\to\infty$,
 $c\nu_{\mathrm{eff}}+\bar B\ge ac\nu_*$ for a fixed $a>0$.
 Along the transport characteristic of the output momentum, the
 integrating factor bounds the source contribution by $C/(c\nu_*)$.
 Boundedness of $c\nu_*q_c(0)$ gives the same bound for the initial
 contribution. For the gain term, the square-row estimate and the
 spatial embedding for $H_X^3L_k^2$ give
 \[
 |K_{\mathrm{eff}}h(s,X,k)|
       \le C\nu_*(k)^{1/4}\|h(s)\|_{H_X^3L_k^2}.
 \]
 Integrating this against the exponential loss gives
 $C\nu_*(k)^{-3/4}\sup_s\|h(s)\|_{H_X^3L_k^2}$.
 Equation~\eqref{hyd:enhanced-bound} bounds the final norm by
 $O(c^{-3/4})$, proving the claim.
\end{proof}

\subsection{Continuation on the prescribed interval and convergence rates}
\label{hyd:completion-section}

\begin{proof}[Proof of Theorem~\ref{hyd:main}]
 Choose all parameters in the following order. First fix the prescribed
 Euler interval and its compact positive parameter neighborhood. Let
 $N_\pm,F_\pm,A_*,B_*$ denote the known bounds in
 Lemma~\ref{hyd:corner}. Set
 \begin{equation}\label{hyd:ratio-choice}
 d_- =\min\{1/4,F_-A_*/(4N_+)\},\qquad
 d_+ =\max\{4,2F_+(B_*+1)/N_-\}.
 \end{equation}
 These constants are fixed before $\Lambda,\kappa_b,c$. With $d_\pm$
 fixed, choose the order weights and geometric thresholds in
 Proposition~\ref{hyd:energy}. Next choose $\Lambda$ large and
 $\kappa_b,e_*$ small enough that
 $C_K(\kappa_b^{1/3}+\Lambda^{-1}+\sqrt{e_*})$ is below half the
 threshold in \eqref{hyd:energy-smallness}, and that the parameter
 $C^1$ error keeps the parameters within the positive neighborhood.
 For the interpolation family relative to $\bar N$ in
 Lemma~\ref{hyd:bulk}, take the uniform thresholds corresponding to
 the slightly wider bounds $d_-/2,2d_+$, and also require
 $C\kappa_b+C_K/\Lambda+C\sqrt{e_*}$ to lie below its smallness
 threshold outside the corner layers. These finitely many requirements
 can be met simultaneously.

 Fix
 \begin{equation}\label{hyd:eta-choice}
 \frac13<\beta<\frac12,\qquad \eta=c^{-\beta}.
 \end{equation}
 For each $c$, this layer is kept fixed throughout the time interval.
 Finally, take $c$ sufficiently large that $c\ge\Lambda$, the geometric
 thresholds hold, $\omega_\eta$ is sufficiently small, and
 $cv_\eta\ge C_K/\kappa_b$. This is possible because
 $cv_\eta\asymp c^{1-2\beta}(\log c)^2\to\infty$.
 The full order of choices is therefore
 \[
 \begin{aligned}
 (R,T,\Theta,K_5)&\longrightarrow d_\pm
              \longrightarrow(\omega_l)_{l=0}^3\\
 &\longrightarrow(\Lambda,\kappa_b,e_*)
              \longrightarrow\beta\longrightarrow c\ge c_*.
 \end{aligned}
 \]
 In particular, $\eta=c^{-\beta}$ is determined last. Both $\Lambda$
 and the factor $\Lambda^{3/2}$ in the energy estimate remain fixed;
 the prescribed $T$ is not shortened to remove this factor.

 Lemma~\ref{hyd:local} constructs the same maximal solution from the
 original initial data. Initially $N_c=\bar N$, and the inverse moment
 map is defined by \eqref{hyd:moment-injective} and the nonsingular
 Jacobian $-M$. Let $t_*$ be the first time this solution violates
 \eqref{hyd:bootstrap-assumptions} or the parameter-neighborhood
 condition, or reaches its maximal existence endpoint. The initial
 distribution ratio is one, and
 $E(0)=O(\Lambda^{-1/2}c^{-3/2})$; hence all these conditions hold
 strictly at the initial time for large $c$.

 For every $t<t_*$, estimate~\eqref{hyd:enhanced-bound} holds with
 constants depending only on the prescribed $T$. By
 \eqref{hyd:eta-choice},
 \begin{align*}
 c\eta v_\eta&=O(c^{1-3\beta}(\log c)^2)\longrightarrow0,\\
 c^{-1/4}v_\eta^{-1/4}
       &=O(c^{-1/4+\beta/2}(\log c)^{-1/2})\longrightarrow0.
 \end{align*}
 Lemma~\ref{hyd:corner} therefore gives
 $r_c^{\mathrm{cor}}\le\min(A_*/2,1)$. By \eqref{hyd:ratio-choice},
 $2d_-\le f_c/N_c\le d_+/2$ on the corner layers.
 Outside those layers, use
 \[
 f_c/N_c-1=(\bar N/N_c)(q_c+h)+(\bar N/N_c-1).
 \]
 Lemma~\ref{hyd:bulk}, the bound $|q_c|\le C/(cv_\eta)$, and the
 parameter energy give
 \begin{equation}\label{hyd:bulk-ratio-improvement}
 \sup_{k\notin E_\eta}|f_c/N_c-1|
 \le C_{K,\Lambda,T}\{(cv_\eta)^{-1}
          +c^{-3/4}v_\eta^{-3/4}+c^{-3/4}\}\longrightarrow0.
 \end{equation}
 The limit uses $-3/4+3\beta/2<0$. Increasing $c$ further makes the
 right-hand side less than $\kappa_b/2$, the energy less than $e_*/2$,
 and the $C^1$ norm of the parameter difference strictly less than
 the margin of the chosen positive neighborhood.

 Thus $t_*<T$ cannot be the first time a temporary inequality becomes
 an equality. The continuity used here is $C_tB_3$ continuity at each
 fixed $c$ and continuity of the local inverse moment map, both supplied
 by the local construction. If the parameter branch were to reach
 the boundary first, the error estimate and $H^3\hookrightarrow C^1$
 would keep its range in a fixed compact subset of the interior of
 $\Theta$. The original inverse function theorem would then continue
 the same branch at that endpoint, excluding this type of exit as well.
 If $t_*$ is the maximal existence endpoint, the bounds just obtained give
 \[
 2d_-N_-\le f_c\le(d_+/2)N_+.
 \]
 Lemma~\ref{hyd:local} then continues the same solution and its finite
 spatial derivatives. The parameters remain in a compact positive set,
 so the inverse moment map extends simultaneously along the same branch.
 This also excludes the maximal endpoint and proves existence through
 $T$. Every comparison function and residual integral is defined from
 the original initial time.

 With $\Lambda,T$ now fixed, \eqref{hyd:enhanced-bound} gives
 $\|e\|_{L^\infty_tH_X^3L_k^2}+\|a\|_{L^\infty_tH_X^3}
 \le C_Tc^{-3/4}$ and
 $\|e\|_{L^2_tH_X^3H_\nu}\le C_Tc^{-5/4}$.
 Combining these with \eqref{hyd:G-integral} proves
 \eqref{hyd:main-H3} and \eqref{hyd:main-weighted}.

 It remains to establish the sharper macroscopic $H^2$ rate. The
 integrability $\nu_*^{-1}\in L^1$ and the moment integrals defining
 the remainder flux give
 \begin{equation}\label{hyd:Z-bounds}
 \|Z^e\|_{L^1_tH_X^3}\le C_Tc^{-5/4},\qquad
 \|Z^G\|_{L^1_tH_X^3}\le C_Tc^{-1}.
 \end{equation}
 Subtracting the Euler equation yields
 \[
 a_t+\sum_jA_j(\theta)\partial_ja
 =-\sum_j[A_j(\theta)-A_j(\bar\theta)]\partial_j\bar\theta
                           +M^{-1}\sum_j\partial_jZ_j.
 \]
 Use the second-order energy with weight $M(\theta)$,
 $\mathcal M_2=\sum_{|\alpha|\le2}
 \int(\partial^\alpha a)^{\mathsf T}M\partial^\alpha a\dd X$.
 Uniform bounds for $M_t$ and $J_{j,X}$ require only
 $\theta_t\in H^2$ and $\theta\in H^3$. The worst commutator is
 $\partial^2 A\,\partial a$, estimated in $L^6\times L^3$.
 The $H^2$ norm of $M^{-1}\operatorname{div}Z$ is controlled by
 $C\|Z\|_{H^3}$. Hence
 \[
 \mathcal M_2'\le C\mathcal M_2
                   +C\sqrt{\mathcal M_2}\|Z\|_{H^3}.
 \]
 Integrate the inequality for $\sqrt{\mathcal M_2+\delta^2}$ and let
 $\delta\downarrow0$. Using $a(0)=0$ and \eqref{hyd:Z-bounds}, we obtain
 $\|a\|_{L^\infty_tH^2}\le C_T/c$, which is \eqref{hyd:main-H2}.
\end{proof}

\begin{corollary}[Entropy and dissipation]\label{hyd:entropy}
 The solution in Theorem~\ref{hyd:main} satisfies
 \begin{align}
 \sup_{t\le T}\int_{\T^3\times D}
     [f_c/\bar N-1-\log(f_c/\bar N)]\,\dd X\dd k
       &\le C_Tc^{-3/2},\label{hyd:entropy-bound}\\
 \int_0^T\frac14\int_{\T^3}\int
     \prod_{i=0}^3 f_{c,i}\,[\Delta(1/f_c)]^2
                   \,\dd\Xi_D\dd X\dd t
       &\le C_Tc^{-2}.\label{hyd:dissipation-bound}
 \end{align}
\end{corollary}
\begin{proof}
 Uniform positive two-sided bounds make the entropy density comparable
 to $|f_c-\bar N|^2$, so the first estimate follows from
 \eqref{hyd:main-H3}. Set $r=(f_c-N_c)/N_c$ and $w=r/(1+r)$.
 Since $\Delta(1/N_c)=0$,
 $\Delta(1/f_c)=-\Delta(w/N_c)$.
 The fixed two-sided comparison of $\prod f_{c,i}$ with
 $\prod N_{c,i}$, boundedness of $\mathfrak a_N$, and $|w|\le C|r|$ give
 \[
 \frac14\int\prod f_{c,i}[\Delta(1/f_c)]^2\,\dd\Xi_D
     \le C\|f_c-N_c\|_{H_\nu}^2.
 \]
 The second estimate follows from \eqref{hyd:main-weighted}.
\end{proof}

\begin{remark}\label{hyd:corner-trace-remark}
 The convergence topology in the theorem allows the corner trace to persist:
 \[
  f_c(t,X,\xi)=N_{\theta_0(X-2\xi t)}(\xi).
 \]
 This trace generally differs from $\bar N(t,X,\xi)$. The corner-layer
 estimates provide positive two-sided comparison bounds and do not
 require the trace to vary with $c$. The time interval of the
 macroscopic result is determined by the smooth positive lifespan
 of the prescribed Euler solution.
\end{remark}
